\begin{trapbox}
  \textbf{\textcolor{CTrap}{易错 1：把不能组成三角形的三个数硬代入。}}

\begin{center}
\begin{tikzpicture}[scale=.9,font=\small,line join=round]
\coordinate (A) at (0,0);\coordinate (B) at (5,0);\coordinate (C) at (3.2,2.4);
\fill[blue!10] (A)--(B)--(C)--cycle;
\draw[thick] (A)--node[midway,below=5pt] {$5$} (B)
 --node[midway,right=5pt] {$3$} (C)--node[midway,left=5pt] {$4$} cycle;
\node[anchor=west,align=left] at (6.1,1.1) {$3+4>5$\\[4pt]能围成三角形};
\draw[very thick,CStatement] (0,-2)--node[midway,above=6pt] {$2$} (2,-2);
\draw[very thick,CExample] (2,-2)--node[midway,above=6pt] {$3$} (5,-2);
\foreach \x in {0,2,5}\fill (\x,-2) circle[radius=1.5pt];
\node[below=6pt] at (2.5,-2) {第三边长 $5$};
\node[anchor=west,align=left] at (6.1,-2) {$2+3=5$\\[4pt]只能排成直线};
\draw[very thick,CStatement] (0,-4.4)--node[midway,above=6pt] {$3$} (3,-4.4);
\draw[very thick,CExample] (3,-4.4)--node[midway,above=6pt] {$4$} (7,-4.4);
\draw[very thick,CTrap,dashed] (7,-4.4)--node[midway,above=6pt] {差 $1$} (8,-4.4);
\draw[thin] (0,-5)--(8,-5);
\node[below=6pt] at (4,-5) {第三边长 $8$};
\node[anchor=west,align=left] at (0,-6.4) {$3+4<8$：另外两段即使拉直，仍够不到第三边的两端。};
\end{tikzpicture}
\end{center}

  如图，$3+4<8$ 时，两段即使拉直也不够长，不能围成三角形。
  公式不是判断三边能否组成三角形的替代品；\textbf{先检查最小两边之和是否大于最大边}。
  特别地，$2,3,5$ 虽然能在公式中算出 $S=0$，但这三段只能首尾排成直线，\textbf{不是真正的三角形}。

  \medskip
  \textbf{\textcolor{CTrap}{易错 2：把半周长当作周长。}}

  $p=(a+b+c)/2$，不是 $a+b+c$。计算时先单独写一行 $p=\cdots$，后面再算 $p-a,p-b,p-c$。

  \medskip
  \textbf{\textcolor{CTrap}{易错 3：漏根号，或把面积取负值。}}

  是 $S^2=p(p-a)(p-b)(p-c)$，因此
  \[
    S=\sqrt{p(p-a)(p-b)(p-c)}>0.
  \]
  由于是\textbf{面积}，不取负根。

  \medskip
  \textbf{\textcolor{CTrap}{易错 4：误以为只要题目出现“三角形面积”，海伦就一定最快。}}

  直角三角形知道两条直角边：直接 $S=\tfrac12 ab$。
  知道两边与夹角：直接 $S=\tfrac12 ab\sin C$。
  海伦公式真正适合的是\textbf{三边已知，角或高反而不方便}。

  \medskip
  \textbf{\textcolor{CTrap}{易错 5：把半周长的差写错。}}

  例如 $p-a=\dfrac{b+c-a}{2}$；它不是 $\dfrac{a+b-c}{2}$。
  可用 $a=y+z,b=z+x,c=x+y$ 重新检查：$p-a=x$，其余同理。

  \medskip
  \textbf{\textcolor{CTrap}{易错 6：求内切圆半径时，和外接圆半径混淆。}}

  $r=S/p$ 说的是\textbf{内切圆}，来自 $S=\tfrac12ar+\tfrac12br+\tfrac12cr$。
  它不能随便拿来当作外接圆半径公式。

  \medskip
  \textbf{\textcolor{CTrap}{易错 7：反求边长时，只解方程，不检查有没有三角形。}}

  面积方程有时会产生不合条件的根。必须把解代回\textbf{三边为正、两边之和大于第三边、周长及面积符合题意}三项条件中核对。

  \medskip
  \textbf{\textcolor{CTrap}{易错 8：记住公式却不理解平方差从哪里来。}}

  在证明中出现
  \[
   4a^2b^2-(a^2+b^2-c^2)^2
  \]
  不要全部展开成四次多项式。直接看成 $u^2-v^2$，先变成两个括号，再各做一次平方差，才会得到四个边长和差。

  \medskip
  \textbf{\textcolor{CTrap}{易错 9：把证明二误认为初中方法。}}

  证明二包含 ``正切两角和公式''，属于\textbf{高中三角函数}；初中阶段想独立证明海伦公式，请使用\textbf{证明一：勾股定理＋平方差}。
\end{trapbox}
